\section{Problem Formulation}
\label{sec:problem}

\subsection{Dynamical System Modeling via ODEs}
Dynamical system modeling provides a mathematical framework for analyzing time-varying phenomena
across the physical, engineering, and biological sciences. Such systems are typically governed by
a system of ordinary differential equations (ODEs),
\begin{equation}
\label{eqn:ode}
    \dot{\vx}(t) = f\big(\vx(t); \vtheta\big), \qquad
    \vx(t) = \begin{bmatrix} \vx_{\mathrm{obs}}(t)\\ \vx_{\mathrm{hid}}(t)\end{bmatrix} \in \mathbb{R}^{d},
\end{equation}
where $\vx(t)$ is the state vector at time $t$, $\vtheta \in \mathbb{R}^{p}$ denotes the
time-invariant parameters governing the dynamics, and $f$ is a sufficiently smooth vector field
ensuring existence and uniqueness of the solution trajectory.

In classical system identification, one estimates the true parameters $\vtheta$ by observing the
state trajectory $\vx(t)$ over an interval $[0,T]$. Once $\vtheta$ is identified, the physical
properties of the system can be analyzed and simulated. In practice, however, full and continuous
measurement of $\vx(t)$ is rarely available. We formalize this through two distinct bottlenecks.

\paragraph{Partial Observability (Non-identifiability).}
The full state is rarely observable in its entirety; it is partitioned as
\begin{equation}
    \vx(t) = \begin{bmatrix} \vx_{\mathrm{obs}}(t) \\ \vx_{\mathrm{hid}}(t) \end{bmatrix},
    \quad \vx_{\mathrm{obs}}(t) \in \mathbb{R}^{d_{\mathrm{obs}}},
    \quad \vx_{\mathrm{hid}}(t) \in \mathbb{R}^{d_{\mathrm{hid}}},
    \quad d = d_{\mathrm{obs}} + d_{\mathrm{hid}},
\end{equation}
where $\vx_{\mathrm{obs}}$ collects the observable states and $\vx_{\mathrm{hid}}$ the unobserved
(hidden) ones. For instance, in the Oral Glucose Tolerance Test (OGTT) a physiological model
couples blood glucose and insulin ($\vx = [x_{\mathrm{glucose}}, x_{\mathrm{insulin}}]^\top$), yet
clinical protocols often measure glucose only ($\vx_{\mathrm{obs}} = x_{\mathrm{glucose}}$). The
absence of $\vx_{\mathrm{hid}}$ induces severe \textit{non-identifiability}: infinitely many
configurations of $\vtheta$ can yield an identical observable trajectory, preventing deterministic
optimizers from converging to the true parameter.

\paragraph{Sparse Sampling (Temporal Aliasing).}
Continuous monitoring is likewise impractical; states are collected only at discrete, sparse times
$\mathcal{T} = \{t_0, t_1, \dots, t_N\}$ with $0 = t_0 < \dots < t_N = T$, yielding
\begin{equation}
    \mX_{\mathrm{obs}} = \big[\, \vx_{\mathrm{obs}}(t_0),\, \dots,\, \vx_{\mathrm{obs}}(t_N) \,\big]
    \in \mathbb{R}^{d_{\mathrm{obs}} \times (N+1)}.
\end{equation}
Sparse sampling introduces temporal aliasing, making rapidly oscillating trajectories
indistinguishable. Moreover, classical pipelines that rely on numerical differentiation
($\dot{\vx}(t) \approx [\vx(t+\Delta t) - \vx(t)]/\Delta t$) incur devastating approximation errors,
rendering gradient-based dynamics learning highly unstable.

\subsection{Goal: Joint Posterior Estimation via Coupled Generative Models}
To estimate $\vtheta$ under partial and sparse observation, we recast parameter estimation as
probabilistic inference: our goal is to characterize the posterior
$p(\vtheta \mid \mX_{\mathrm{obs}})$ given only the sparse, partial data $\mX_{\mathrm{obs}}$.

Because $\mX_{\mathrm{hid}}$ is entirely absent, directly optimizing
$p(\vtheta \mid \mX_{\mathrm{obs}})$ is ill-posed and numerically intractable. We therefore model
the \textit{joint posterior} over both the parameters and the missing hidden trajectory,
$p(\vtheta, \mX_{\mathrm{hid}} \mid \mX_{\mathrm{obs}})$, where
\begin{equation}
    \mX_{\mathrm{hid}} = \big[\, \vx_{\mathrm{hid}}(t_0),\, \dots,\, \vx_{\mathrm{hid}}(t_N) \,\big]
    \in \mathbb{R}^{d_{\mathrm{hid}} \times (N+1)}.
\end{equation}
We decompose this joint estimation into two coupled conditional generative models, parameterized by
$\phi$ and $\psi$:
\begin{align}
    \text{\textbf{Hidden Trajectory Generator:}} \quad & p_\phi(\mX_{\mathrm{hid}} \mid \mX_{\mathrm{obs}}, \vtheta), \\
    \text{\textbf{Parameter Generator:}} \quad & p_\psi(\vtheta \mid \mX_{\mathrm{obs}}, \mX_{\mathrm{hid}}).
\end{align}
Given an initial estimate $\vtheta^{(0)}$, we approximate the intractable joint posterior by an
alternating sampling scheme, for $k = 0, 1, 2, \dots$:
\begin{align}
    \mX_{\mathrm{hid}}^{(k+1)} & \sim p_\phi(\mX_{\mathrm{hid}} \mid \mX_{\mathrm{obs}}, \vtheta^{(k)}), \label{eq:update_x} \\
    \vtheta^{(k+1)} & \sim p_\psi(\vtheta \mid \mX_{\mathrm{obs}}, \mX_{\mathrm{hid}}^{(k+1)}). \label{eq:update_theta}
\end{align}
This process refines the uncertainty in both the hidden-state and parameter spaces simultaneously,
guiding the system toward feasible parameters rather than a single point estimate.

\paragraph{Amortized Inference via Synthetic Simulation.}
To avoid per-instance MCMC or online grid search, we adopt an \textit{amortized} framework. We
assume access to a synthetic dataset
$\mathcal{D} = \{ (\mX_{\mathrm{obs}}^{(m)}, \mX_{\mathrm{hid}}^{(m)}, \vtheta^{(m)}) \}_{m=1}^{M}$,
generated by sampling parameters $\vtheta^{(m)}$ and initial conditions $\vx^{(m)}(0)$ from a
physically plausible prior $p(\vtheta, \vx_0)$ and integrating the forward ODE,
\begin{equation}
    \vx^{(m)}(t) = \vx^{(m)}(0) + \int_{0}^{t} f\big(\vx^{(m)}(\tau); \vtheta^{(m)}\big)\, d\tau.
\end{equation}
Training $p_\phi$ and $p_\psi$ on $\mathcal{D}$ amortizes the inference cost into an offline phase;
once trained, the networks approximate the joint posterior and enable instantaneous,
distribution-aware recovery for any newly observed sparse sequence $\mX_{\mathrm{obs}}$.
